\section{Protocol}\label{sec:protocol}

This section sets out the measurement protocol. Each rule addresses a
confound that we observed on this family, and each is stated together with the
evidence that motivates it.

\subsection{Native stopping rules as a confound}\label{sec:native}
Suppose every solver is run to its own exit criterion and the resulting costs
are compared. The comparison then credits a method for stopping early as much
as for converging quickly. This happens even when all methods share one
nominal rule. Under the dimension-scaled rule $\nrm{\nabla\theta}_2\le10^{-7}n$
used by~\cite{huynh2025accelerated}, which at $n=100$ is $10^{-5}$, the six
solvers of \cref{sec:ranking} and the Newton-SIN control stop at median
forward errors between
$5.04\times10^{-8}$ and $1.93\times10^{-5}$, a factor of 383
(\cref{sec:stopping}, \cref{fig:native}). A cost comparison taken at these
exits compares different accuracies.

\subsection{Work--accuracy trajectories against the exact solution}\label{sec:traj}
We take one \EVD\ as the unit of cost. After \emph{every} \EVD, whether it
produces an iterate, a line-search trial or a Jacobian factorization, we record
the pair
\[
  \bigl(\text{cumulative \EVDs},\ \fro{X_k-\Xs}\bigr),
\]
together with $\nrm{\nabla\theta}_2$ and a flag indicating whether the
evaluated point was accepted. Here $\Xs$ is the exact solution from
\cref{prop:exact}, not the output of another solver. Measuring forward error
rather than a residual follows the standard distinction
of~\cite[Ch.~1--2]{higham2002accuracy}. All runs use a common,
tight tolerance of $10^{-11}$, and each solver has its own cap (\cref{sec:ranking}). These trajectories are work--precision diagrams in the sense
of~\cite{hairer1996solving}, and reading costs off them at a ladder of targets
follows data profiles~\cite{more2009benchmarking} and fixed-target
evaluation~\cite{hansen2021coco}. The difference is that the reference is the
exact solution, not the best value found by any solver. For any target
$\varepsilon$, the cost can then be read off the trajectory after the run. We
define the \emph{cost to $\varepsilon$} as the cumulative \EVD\ count at the
first accepted iterate from which $\fro{X_k-\Xs}\le\varepsilon$ holds for
every later accepted iterate, that is, at the start of the final sub-$\varepsilon$
stretch of the run. A solver whose last accepted iterate lies above the target
is recorded as not reaching it at that $\varepsilon$. A run that dips below the
target, leaves it, and settles below it again is credited from the point at
which it settles.

Trajectories also expose \emph{plateaus}. The cost of a quadratically
convergent method takes very few distinct values across many decades of
$\varepsilon$: in \cref{tab:fwd}, Newton-SIN-BH costs 3, 4 or 5 \EVDs\ at
every target from $10^{-2}$ to $10^{-10}$. A single-tolerance comparison
therefore depends on where the tolerance falls relative to these plateaus.

\subsection{Accepted-only accounting}\label{sec:accounting}
A rejected line-search trial always costs one \EVD. Under
\emph{accepted-only} accounting it can never be credited with reaching a
target, because the method did not move to that point. Under
\emph{all-trial} accounting it can, which credits a method for an accuracy it
never returned. We use accepted-only accounting throughout and report the
other convention as a sensitivity check (\cref{sec:acct-results}); on the
family studied here the difference turns out to be small.

\subsection{Primitive work as the unit of cost}\label{sec:primitive}
\Cref{tab:primitive} gives the median primitive work at matched tolerance
($10^{-11}$) over the 270 instances.

\begin{table}[t]
\centering
\caption{Median primitive work at matched tolerance $10^{-11}$, 270 instances,
$n=100$. An `outer iteration' denotes a different amount of work in each
row.}
\label{tab:primitive}
\small
\begin{tabular}{lrrrr}
\toprule
solver & \EVDs & CG iterations & line-search trials & accepted outer its\\
\midrule
Newton-SIN-BH & 5     & 24   & 0  & 4\\
Newton-SIN    & 11    & 29.5 & 0  & 5\\
AGD-SDAJ-BH   & 70    & 0    & 14 & 9\\
AGD-SDAJ      & 81.5  & 0    & 25 & 10\\
SBB-Dual      & 46.5  & 0    & 0  & 45.5\\
Dykstra-APM   & 102.5 & 0    & 0  & 102.5\\
Anderson-APM  & 34    & 0    & 0  & 34\\
\bottomrule
\end{tabular}
\end{table}

Four accepted Newton iterations use 5 \EVDs\ (one at the initial point plus one
per accepted step) and 24 Krylov iterations. SBB-Dual,
Dykstra-APM and Anderson-APM use one \EVD\ per iteration and no further
primitive, although Anderson-APM's least-squares update carries $O(mn^2)$ work
that the \EVD\ count does not record. Nine AGD-SDAJ-BH outer
iterations use 70 \EVDs, because each outer iteration contains $q=2$
quasi-Newton inner steps and one accelerated gradient step, each with its own
line search, and the accelerated step can need a further evaluation
(\cref{alg:agd}). A table of
`iterations to convergence' across these rows would compare a different
quantity in each method class (Newton, AGD, SBB, APM). We therefore report \EVDs, Krylov iterations and
line-search trials separately.

\subsection{Timing protocol}\label{sec:timing}
All timings use Julia with \texttt{BLAS.set\_num\_threads(1)}. A warm-up solve
is run and discarded, \texttt{GC.gc()} is called before the measured solve,
and a single \texttt{@elapsed} surrounds the whole solve. A separate
instrumented pass (TimerOutputs) gives the breakdown of time. In that
breakdown, percentages are taken against the measured elapsed time, so any
uninstrumented work shows up as `other' time instead of inflating the
\EVD\ share. Timings were measured on all 270 instances at $n=100$, and on an
18-instance subset at $n=500$ (listed in \texttt{code/n500\_timing\_subset.txt}),
at both the native ($10^{-7}n$) and the matched ($10^{-11}$) tolerance. The
trajectory runs are a separate, untimed pass.

All seven variants were timed back to back in a single session, because
timings from different sessions are not comparable: an earlier session on the
same machine, with identical \EVD\ counts on every run, measured the dual
methods, whose code had not changed, between 3.7 and 4.0 times slower per
\EVD\ at $n=100$. Each
method is timed in an implementation that does not add work of its own. The
projection methods rebuild the projection from the positive eigenpairs, as the
dual methods do, rather than from all $n$, which removes an $O(n^3)$ product
per \EVD\ and gives bit-identical iterates on all 1104 projection-method runs
checked (270 $n=100$ instances and 6 $n=500$ instances, $\times$ 2 methods
$\times$ 2 tolerances; \texttt{code/validation/check\_projection\_rebuild.jl}); and
Anderson-APM is timed in a form that updates its buffers in place, with
iterates bit-identical to the transcription verified in \cref{app:anderson}.
Timed as transcribed, it would have cost $2.3$ to $3.8$ times more per \EVD.
All timings were measured on one machine: Intel Core 7 150U (10 physical
cores, 12 logical), 32\,GB RAM, Windows 11 (build 26200), Julia 1.12.6 with
OpenBLAS (ILP64) through libblastrampoline. The timing runs pin
\texttt{BLAS.set\_num\_threads(1)}; the untimed trajectory pass at $n=500$ was
run with four BLAS threads, which changes wall-clock only and not the \EVD\
counts it records.

\subsection{Fairness of the \EVD\ as a unit of work}\label{sec:evd-unit}
The \EVD\ count is a good proxy for cost only if one \EVD\ costs about the same
in every method. \Cref{tab:perevd} and \cref{fig:perevd} give the median
wall-clock time per \EVD\ at matched tolerance.

\begin{table}[t]
\centering
\caption{Microseconds per \EVD\ at matched tolerance ($10^{-11}$), $n=100$,
single-threaded BLAS.}
\label{tab:perevd}
\small
\begin{tabular}{lrr}
\toprule
solver & \si{\micro\second}/\EVD & relative to cheapest\\
\midrule
SBB-Dual      & 1133 & 1.00\\
AGD-SDAJ-BH   & 1181 & 1.04\\
AGD-SDAJ      & 1188 & 1.05\\
Dykstra-APM   & 1278 & 1.13\\
Anderson-APM  & 1323 & 1.17\\
Newton-SIN    & 1341 & 1.18\\
Newton-SIN-BH & 1695 & 1.50\\
\bottomrule
\end{tabular}
\end{table}

\begin{figure}[t]
\centering
\begin{tikzpicture}
\begin{axis}[
  width=0.8\linewidth, height=5.6cm,
  xbar, bar width=9pt,
  xmin=0, xmax=2000,
  xlabel={\si{\micro\second} per \EVD\ (matched tolerance $10^{-11}$)},
  ytick=data,
  % labels come from the data file, so they cannot fall out of step with the bars
  yticklabels from table={figures/data/per_evd_cost.dat}{name},
  y dir=reverse,
  nodes near coords, nodes near coords align={horizontal},
  every node near coord/.append style={font=\footnotesize},
  enlarge y limits=0.1,
  axis x line*=bottom, axis y line*=left,
]
\addplot[fill=gray!55, draw=gray!80!black] table[x=usperevd, y=idx] {figures/data/per_evd_cost.dat};
\end{axis}
\end{tikzpicture}
\caption{Median wall-clock time per \EVD\ at $n=100$, all seven variants timed
in one session. The spread is a factor of 1.50; Newton's \EVDs\ are the
dearest, because they are accompanied by Krylov solves that the \EVD\ count
does not record. Anderson-APM's \EVDs\ cost 1.17 times SBB-Dual's, the
difference coming from its work outside the eigendecomposition, chiefly its
least-squares update of $O(mn^2)$ work per \EVD.}
\label{fig:perevd}
\end{figure}


The median is taken over every run, including any that exhausted its budget
without converging (notably Newton-SIN's 56 naive-Armijo stalls of \cref{sec:trap2}).
In these timing tables a run's cost to its own exit counts as its cost ---
unlike the cost-to-target rule of \cref{sec:traj}, under which
a run that never reaches a target is recorded as not reaching it.
\Cref{tab:perevd} shows a factor-of-1.50 spread in per-\EVD\ cost across the
seven variants. Two solvers carry work the \EVD\ count does not see, and in
both cases counting \EVDs\ overstates that solver's advantage without
reversing it: Newton-SIN-BH is cheaper than SBB-Dual by 9.3 in \EVDs\
(5 against 46.5) but only 5.4 in wall-clock time (0.009\,s against
0.050\,s), because each Newton outer iteration hides an uncounted Krylov
solve; and Anderson-APM, whose least-squares update costs about 13--15\% of
its time against about 1\% for SBB-Dual, uses about 27\% fewer \EVDs\ than
SBB-Dual but only about 5\% less wall-clock time. In this single session, then,
the known bias in \EVD-counting is smaller than the gaps it is used to
compare. We report both metrics throughout, but for the reason given next we
treat \EVDs\ as primary and wall-clock time as a secondary check rather than
the reverse.

\paragraph{Session-to-session variation.}
\Cref{tab:perevd} is a single measured session. Repeating the
matched-tolerance timing measurement four further times shows a given
solver's median microseconds/\EVD\ varies far more across sessions (a
factor of 3.8 to 4.9, depending on the solver and dimension;
\texttt{results/analysis/timing\_repeats.md}) than the $\le\!1.50\times$
spread between solvers within one session --- a swing large enough, and
uniform enough across solvers within a session, that we read it as a
platform effect (\cref{sec:discussion}) rather than a property of the
methods. Two features nonetheless survive it: the cheapest solver per \EVD\
is always SBB-Dual or an AGD variant, at both dimensions, and the dearest per
\EVD\ is Newton-SIN-BH in 8 of the 10 sessions (Anderson-APM and Dykstra-APM
each take it once). It is the middle of the ranking that reorders, so the
specific ratios in \cref{tab:perevd} should be read as one session's draw,
not stable constants; the 5.4-in-time figure above is likewise approximate
to about ten per cent across sessions, while the
9.3-in-\EVDs\ figure, involving no timing measurement, is unaffected.
Ranking the seven variants by median \EVDs\ and by median wall-clock time
agrees in only 5 of these 10 sessions. The ranking by median \EVDs\ is
identical in all 10 sessions, since \EVD\ counts do not vary across sessions,
so every disagreement is wall-clock noise. By this second, wall-clock notion of
``dearest,'' Newton-SIN-BH is still cheapest overall in every session, but
the dearest \emph{overall} solver is Dykstra-APM in 8 of the 10 (a different
8 of 10 from the per-\EVD\ count above: all five at $n=500$, three of five
at $n=100$, where AGD-SDAJ overtakes it twice despite needing about a fifth
fewer \EVDs). So only the identity of the single cheapest method is
independent of the choice of cost metric; even within one session the
metric can decide a close pair (at $n=500$, Anderson-APM needs fewer \EVDs\
than AGD-SDAJ-BH on 13 of 18 instances but less time on only 11).

\subsection{Summary of the protocol}
\begin{enumerate}
  \item Measure forward error against an exact $\Xs$.
  \item Run every solver to a common tight tolerance and record a trajectory
    point after every \EVD.
  \item Read costs off the trajectories at a common target; give rejected
    trials no accuracy credit.
  \item Report \EVDs, Krylov iterations, line-search trials and wall-clock
    time, never iterations alone.
  \item Report the feasibility of the returned iterate
    ($\lambda_{\min}(X)$ and $\nrm{\diag(X)-e}_\infty$), and state which
    iterate is returned.
  \item Report reach rates beside every median that is computed over a
    subset of instances.
\end{enumerate}
